\section{Local clocks, record gluing and information capacity}
\label{app:record-gluing-principle}

The proposed record-gluing law RG connects response directions to classical
record transport. It has distinct information, geometry and implementation
premises. A simpler sufficient transport principle replaces exact dilation
covariance and continuous rotation symmetry by a regular local clock and
time-faithful refinement. This is a proposed physical extension, not a fourth
adopted core axiom or a property certified for the simulator. Its target is
the operational causal/count limit; it does not uniquely select a finite
Fibonacci population or a microscopic read radius.

\paragraph{The local-clock interface.}
Let $E$ be the three-dimensional simple response ideal with its positive
trace metric. Its attachment to operational positions is supplied. Assume:
\begin{enumerate}
\item Independent classical bits, clean auxiliary registers, retained
versions, preparation, discard, waiting and NOT/Toffoli gates admit
arbitrary finite composition and parallel realization. All auxiliary work
is counted. Local processing has zero ideal duration or arbitrarily small
positive overhead for each complete finite program.
\item Every sufficiently small displacement $v$ has a retained-record
primitive of duration $\ell(v)$, independent of origin and covariant under
carrier rotations, with $\ell(0)=0$ and
\[
  \ell(v)^2=Q(v)+o(\|v\|^2),\qquad Q>0.
\]
The expansion is uniform in direction. Local directional coverage is a
premise; finitely many rotated directions do not imply it.
\item Every native dependency, including control, timing and metadata,
admits refinement into small primitives and waits with the same endpoints
and elapsed time converging to the original time. A local primitive costs
at least $\ell(v)$. Finite programs have exact value-preserving realizations
with arbitrarily small whole-program overhead. Local interventions at every
designated intermediate output before publication are preserved while all
other existing records and caches are fixed.
\end{enumerate}
These are LC1--LC3. They constrain all admitted influences, including those
outside a chosen observation program.

\paragraph{Finite symmetry and the limiting cone.}
In an orthonormal icosahedral frame, invariance of a symmetric quadratic form
under $\operatorname{diag}(1,-1,-1)$ and
$\operatorname{diag}(-1,1,-1)$ makes its off-diagonal entries zero. The cyclic
coordinate permutation equates its diagonal entries. These proper rotations
preserve the twelve ports and generate a tetrahedral subgroup. Thus
$Q(v)=\tau^2\|v\|^2$ for $\tau>0$, and uniformly
$\ell(v)=\tau\|v\|+o(\|v\|)$.

For any native path of elapsed time $T$ and displacement $v$, refine until
every step satisfies the local lower bound. For any $0<\epsilon<\tau$,
\[
 T_n\geq\sum_j\ell(v_j)
 \geq(\tau-\epsilon)\sum_j\|v_j\|
 \geq(\tau-\epsilon)\|v\|.
\]
Time-faithful refinement and then $\epsilon\to0$ give $T\geq\tau\|v\|$.
Conversely, $N$ equal steps have total time $N\ell(v/N)\to\tau\|v\|$.
Finite realization with a sufficiently small overhead meets every strict
deadline $T>\tau\|v\|$; waiting fills the remaining time. The resulting
access cone has speed $c=1/\tau$. Exact null access is not asserted.

The example $\ell(v)=\tau r+\kappa r^2$, $r=\|v\|$, $\kappa>0$, has
subdivision time $\tau r+\kappa r^2/N$. Subdivision followed by waiting
preserves an old primitive's duration, yet no nonzero finite path attains
the null boundary. It violates exact dilation covariance while satisfying
the local-clock principle. This separates LC from exact RG geometry; an
approximate implementation of RG can have the same limiting behavior.
The icosahedrally invariant polyhedral norm
$\max_{u\in V}|u\cdot v|$ shows why finite symmetry without quadratic
regularity is insufficient. A fast coarse shortcut preserves local Taylor
data but violates time-faithful refinement. Both premises bear physical
content.

For $q^3$ distinguished records in a window of side $L$, spacing $L/q$,
and layers $\delta_q\to0$ with $q\delta_q\to\infty$, finite independent
programming realizes inner read menus
$a_q=(1-\eta_q)c\delta_q$, $\eta_q\to0$, with sufficient deadline slack.
The source-net covering theorem gives inner and outer radii
$k(a_q-2h_q)$ and $ka_q$, where $h_q=\sqrt3L/q$. Localized intermediate
interventions place consumed-writer paths in complete native ancestry;
the clock theorem excludes native spacelike shortcuts. The orders converge
away from the null boundary. Cell assignment supplies count-volume weight
$(L/q)^3\delta_q$ for the distinguished records. The M1 choice
$\delta_q=L/(c\sqrt q)$ is one regulator; auxiliary work has a separate bill.

\paragraph{Classical compilation and finite algebra coherence.}
Every Boolean map has a clean reversible implementation: toggle an output
for each enabled minterm, compute its multi-input conjunction using a
Toffoli ladder, and reverse the ladder. NOT gates implement negative
controls. CNOT is a derived Toffoli gate with a constant control; its clean
control register, preparation and reset are counted. The input is retained, the desired value is XORed into the output,
and all auxiliary bits return to zero. This construction is exponential in
general. It proves finite programmability, not an efficiency or energy bound.

On a finite graph take private factors $\mathbb C^{12}_i\otimes Q_i$ and
shared edge buffers $B_e=M_{12}$. The carrier algebra consisting of its
private factors and incident buffers has center exactly $\mathbb C^{12}_i$.
The channel with Kraus operators
\[
 K_{p,b}=e_p\otimes|p\rangle\langle b|,
 \qquad \sum_{p,b}K_{p,b}^{\dagger}K_{p,b}=I,
\]
retains central port $p$ and copies it into a quantum buffer's fixed basis.
Reading that basis and preparing a target port completes a classical copy;
unknown quantum states are not cloned. Off-diagonal buffer effects remain
usable measurements. Private response actions commute with these copying
maps. Buffer access, controller transport and readback receive their actual
local-clock grades.

An explicitly included joint observer and its marginal constraints provide
a nonempty finite state family. Faithful tracial references and an injective
positive-weight entropy cover give the joint tracial state as the unique
unconstrained minimum. Independent preparations are interventions or
constraints, not that unconstrained minimum's selected records. Adding blank
factors and lifting old channels gives value-level refinement. This is a
finite consistency witness. It does not furnish the all-level source-bound
support, degree-one bridge, operational deletion tests or proof that the
private response exhausts all admitted reversible source responses. A
dormant support added to a circuit would not meet those obligations.

\paragraph{Independent motivation and its boundary.}
Reversible classical computation motivates cleaning auxiliary work after
copying an answer. Its relation to a transport clock is an extra premise.
For a tracial matrix reference $\rho_0=I/d$ and injective traceless Hermitian
linear perturbations $X(v)$, direct expansion gives
\[
 D(I/d+X(v)\Vert I/d)=\frac d2\operatorname{Tr}X(v)^2+O(\|v\|^3).
\]
An equivariant perturbation makes this positive susceptibility isotropic by
the same finite-symmetry argument. Identifying $\ell(v)^2$ with a positive
multiple of this entropy cost to second order would supply a regular clock,
but A3 supplies neither that calibration nor its coefficient. Equal-radius
latency ratios and coarse/fine timing comparisons test this proposal without
an M1 fit. Entropy curvature alone is not evidence for a physical clock law.

\paragraph{Information and resource consequences.}
If a receiver decodes $m$ independently intervenable $d$-ary records from a
finite transcript alphabet $\mathcal T$, with all caches fixed, the transcript
map is injective. Therefore $|\mathcal T|\geq d^m$. Payload, timing, control
and metadata all count. Additive accumulation cannot decode arbitrary tuples;
preexisting correlations do not evade fresh independent interventions.

For $q=n^2$ and $n\geq2$, put $m=\lfloor n/2\rfloor$. In the M1 grid,
all offsets in $[-m,m]^3$ lie strictly within the radius $L/\sqrt q$.
At least $(n^2-2m)^3$ interior receivers each consume $(2m+1)^3$ records.
Thus there are at least $n^9/8=q^{9/2}/8$ read incidences per layer and
$n^{10}/8=q^5/8$ in a horizon $L/c$. Each interior receiver needs at least
$(2m+1)^3-1\geq n^3/2$ incoming binary transcript bits per layer under fresh
independent interventions, hence bandwidth at least $cq^2/(2L)$. This excludes
its own locally available record from the communication cut.
Multicast can share transport paths;
the aggregate statement counts output incidences rather than unicast packets.
The rate bound presumes independent new interventions with prior caches fixed.
It does not infer a fresh entropy rate from an unperturbed correlated history.

Cofinal workspace is consequently substantive. A fixed per-receiver bandwidth
cannot support literal dense lossless observation at all refinements under
these conditions. A single central $\mathbb C^{12}$ record has only twelve
alternatives; an unbounded lossless tuple therefore requires a growing
composite observer/register bank. The logical receiver cannot be identified
with one unchanged primitive carrier. The same transport law also permits identity, parity and
other local updates, so it selects no matter equation, energy law or Lorentz
covariance of dynamics. Its consequences beyond M1 are communication,
refinement timing and information-capacity constraints.

The complete analytic argument and controls are in
\ophid{code/rg_principle/DERIVATION.md}.
The compact receipt independently replays 263 finite circuits, including
7,880 localized interventions, with every reference gate and work register
counted. Exact algebra controls cover all 1,728 basis units of the copying
channel. Twelve kernel-checked results and a transitive standard-axiom audit
are in \ophid{Lean/Geometry/RecordGluingPrincipleAxiomAudit.lean}.
Finite controls do not replace the analytic all-size arguments or certify a
native source realization.
